\section{Experimental Setup}
\label{sec:setup}

\paragraph{Model and training.}
The policy is Qwen2.5-VL-3B-Instruct~\cite{qwen25vl2025}, trained with the multi-turn PPO of VAGEN~\cite{vagen2025} on verl~\cite{hybridflow2025} with vLLM~\cite{vllm2023} for generation.
The observation budget is $D = 4$ (at most $5$ turns per episode), each response is limited to $256$ tokens, the observation cost is $\lambda = 0.05$, and the training change rate is $\rho = 0.5$.
All methods share one frozen configuration: reward, optimizer, learning rate, KL coefficient, $64$ episodes per step, and \tbd\ training steps \todo{from configs/rl\_common.yaml once frozen}.
For the consistency term we use $m = 1$ and $\beta = 0.1$.
Each trained method is run with three seeds.

\paragraph{Methods.}
The $2{\times}2$ design of \cref{sec:credit} gives four methods, all with the same reward and rollout budget:
\begin{itemize}
\item \textbf{PPO} (single-world, critic): standard outcome-based multi-turn PPO.
\item \textbf{Paired-GAE} (paired, critic): the same lockstep pairs, with the critic's GAE in each world. \todo{Run if compute allows; otherwise state that this cell is missing.}
\item \textbf{Single-$\Qc$} (single-world, commit-now baseline): $A^{\mathrm{single}}_t$ before the evidence, our adaptation of a same-world look-versus-commit contrast~\cite{seekvln2026}.
\item \textbf{IS-credit} (paired, commit-now baseline): information-set credit, \cref{eq:is}.
\end{itemize}
\textbf{Ours} adds post-evidence consistency (\cref{sec:consistency}) to IS-credit.
The comparisons are: IS-credit vs.\ Single-$\Qc$ for the effect of paired sampling; Single-$\Qc$ vs.\ PPO (and IS-credit vs.\ Paired-GAE) for the effect of the commit-now baseline; and Ours vs.\ IS-credit for the effect of the consistency term.

\paragraph{Training-free references.}
The base model chooses the target in all of them; they differ only in when it may commit.
\emph{Never look} must commit at the first turn.
\emph{Always look} spends the budget on a fixed sweep (left, right, right, left) that returns to $s_0$, and then commits.
\emph{Look twice} performs the first two actions of that sweep.
\emph{Look until seen} keeps sweeping until the anchor is visible in its own world; it uses ground-truth instance masks and is therefore privileged.
\emph{Base model} decides by itself when to look and when to commit.

\paragraph{Metrics.}
Success rate (SR) is the fraction of episodes that select the correct target; observation cost is the number of observation actions per episode.
\emph{Premature commitment} is the fraction of \Look\ episodes whose \commit\ occurs before $\tauseen$, \ie before the agent has seen the anchor in its own world, including episodes that never see it.
\emph{Post-evidence update accuracy} is the accuracy of \Look\ episodes that commit after $\tauseen$, reported separately for \Look$\to$\Flip\ (correctly switching) and \Look$\to$\Stay\ (correctly keeping), together with the number of such episodes.
We also report accuracy on direct \Flip\ and \Stay\ pairs and SR as a function of observation cost.
Following our pre-registered protocol, the primary endpoints are SR and premature commitment at test change rates of $0.5$ and $0.25$ under a matched observation budget, with the constraint that SR at rate $0$ does not drop significantly.
Each test episode uses the edited world with probability equal to the test change rate, with a fixed sampling seed.
Evaluation uses greedy decoding.
For trained methods we average episodes within each seed and report the mean, standard deviation, and $95\%$ confidence interval over seeds; a difference is called significant only if the $95\%$ confidence interval of the Welch difference excludes zero.
Secondary measures are the variance of the pre-evidence credit per independent decision unit during training and the area under the validation SR curve.
